\documentclass[12pt,a4paper]{article}
\usepackage[UTF8,heading=true]{ctex}
\usepackage{geometry}
\usepackage{amsmath,amssymb,amsthm}
\usepackage{graphicx}
\usepackage{booktabs}
\usepackage{longtable}
\usepackage{array}
\usepackage{enumitem}
\usepackage{hyperref}
\usepackage{fancyhdr}
\usepackage{xcolor}
\usepackage{listings}
\usepackage{setspace}
\usepackage{caption}
\usepackage{float}

\ctexset{
  section = {
    name = {,、},
    number = \chinese{section},
    format = {\Large\bfseries\raggedright}
  }
}

\geometry{left=2.5cm,right=2.5cm,top=2.5cm,bottom=2.5cm}
\setlength{\headheight}{15pt}
\setstretch{1.5}

\pagestyle{fancy}
\fancyhf{}
\fancyhead[C]{发明专利技术交底书}
\fancyfoot[C]{\thepage}

\begin{document}

\begin{center}
    {\Large\heiti 中国移动专利申请} \\[1em]
    {\Huge\heiti 技术交底书} \\[1em]
\end{center}

\noindent
\begin{tabular}{|p{3cm}|p{12cm}|}
\hline
公司编号 &  \\
\hline
发明名称 & 噪声自适应变分量子加密（NAV-QE）系统、方法及其设备指纹与硬件绑定密钥应用 \\
\hline
申报单位 & 中国移动通信有限公司研究院 \\
\hline
申报类型 & 发明 \\
\hline
发明人 & 【请填写正式发明人名单，按专利署名顺序确认】 \\
\hline
技术联系人 & 【请填写联系人姓名、部门、电话及邮箱】 \\
\hline
注意事项 & 1．技术联系人应为深入了解本申请提案技术方案的技术人员。 \\
         & 2．请按照集团公司提供的本技术交底书模板逐项填写。 \\
\hline
\end{tabular}

\vspace{1cm}

\section{发明名称}

噪声自适应变分量子加密（NAV-QE）系统、方法及其设备指纹与硬件绑定密钥应用。

\section{技术领域}

本发明涉及量子计算安全、后量子密码、硬件指纹识别、设备认证与物理不可克隆函数技术领域，具体涉及一种基于量子处理器固有噪声特征构建硬件指纹、生成设备绑定密钥、执行硬件证明并识别篡改行为的噪声自适应变分量子加密系统与方法。

\textbf{与量子计算的关联说明：}

随着量子计算能力持续增强，传统 RSA、ECC 等公钥密码体制面临被 Shor 算法破解的潜在风险，“现在存储、未来解密”已成为现实安全威胁。当前主流的后量子密码方案大多从算法层面对抗量子攻击，但普遍缺少与具体量子硬件实体强绑定的物理信任根。

另一方面，NISQ 时代和未来容错量子阶段的硬件天然存在量子退相干、门误差、读出误差和跨比特串扰。这些特征虽然会影响计算精度，却又因材料缺陷、制造差异、封装结构、冷却条件和运行环境的不同而形成稳定但不可复制的个体差异。如果能够对这些噪声特征进行表征、归一化、量化与映射，即可将其由“计算负担”转化为“安全资源”，进而用于量子硬件身份认证、设备绑定密钥派生、持续证明与篡改检测。

本发明可应用于：
\begin{itemize}
    \item 量子云平台中的设备硬件证明与可信接入
    \item 量子控制平面和量子中间件中的可信执行验证
    \item 量子硬件安全模块（qHSM）中的设备绑定密钥生成
    \item 保密通信系统和受监管基础设施中的量子设备身份管理
\end{itemize}

\textbf{关联项目：}\underline{\hspace{8cm}}（待填写）

\section{现有技术的技术方案}

当前与本发明相关的现有技术主要包括以下几类：

\subsection{量子噪声表征与校准技术}

现有量子计算平台通常通过 $T_1$、$T_2$、随机基准测试、读出误差估计和串扰测量等方法获取量子处理器噪声参数，用于校准优化、误差缓解和编译调优。例如，$T_1$ 可通过激发态衰减实验拟合获得：
\[
P_1(t) = A e^{-t/T_1} + B
\]

$T_2$ 可通过 Ramsey 或 Echo 实验得到，例如：
\[
P_0(t) = \frac{1}{2}\left(1 + A e^{-t/T_2}\cos(\Delta \omega t + \phi)\right)
\]

随机基准测试可提取门误差参数，其拟合形式为：
\[
F(m) = A p^m + B
\]
其中，$m$ 为 Clifford 序列长度，$p$ 为衰减参数。

\subsection{经典 PUF 与设备认证技术}

传统物理不可克隆函数主要基于硅制造偏差、SRAM 上电状态、环振荡器频差或传播延迟差等物理特征生成设备身份标识和挑战—响应对。这类技术适用于经典芯片与嵌入式设备，但通常不适配量子处理器的退相干、门误差和串扰等噪声特征。

\subsection{量子随机性与量子 PUF 方向}

量子随机数发生器能够利用量子测量产生高熵比特流，但其重点通常是随机数质量，而非设备身份绑定。部分量子 PUF 研究从量子态不可克隆性出发讨论认证问题，但尚未形成基于实际量子处理器噪声特征的稳定硬件指纹—密钥—篡改检测闭环工程方案。

\section{现有技术的缺点及本申请提案要解决的技术问题}

现有技术主要存在以下缺陷和技术问题：

(1) \textbf{仅支持噪声测量，不支持安全用途闭环}：现有量子噪声表征工作主要服务于校准、误差缓解和性能优化，未将噪声参数直接转化为设备身份、证明令牌和密钥材料。

(2) \textbf{仅支持硬件指纹概念，不支持稳定密钥重构}：经典 PUF 不适配量子处理器，现有量子 PUF 相关研究也普遍缺少针对连续噪声参数的归一化、量化、模糊提取和稳定密钥派生机制。

(3) \textbf{无法有效区分自然漂移与恶意篡改}：量子硬件运行期间天然存在缓慢漂移，若缺乏统计建模和阈值机制，则难以在不过度误报的前提下识别探针插入、电磁干扰、异常温漂等恶意事件。

(4) \textbf{缺少工程化部署路径}：现有技术大多停留在理论方法、实验指标或单点算法层面，尚未给出面向量子云 API、量子中间件、控制平面和 qHSM 的完整产品化结构。

本申请提案旨在解决上述问题，提供一种可部署的 NAV-QE 系统与方法，将量子硬件噪声资源系统性转化为设备指纹、设备绑定密钥和持续证明能力。

\section{本申请提案的技术方案的详细阐述}

本发明提供一种噪声自适应变分量子加密（NAV-QE）系统，至少包括量子处理器模块、参数化电路执行模块、机器学习表征模块、误差映射与指纹生成模块、密钥生成模块、篡改检测模块以及部署接口模块。

\subsection{系统架构}

其中，量子处理器模块包括多个物理量子比特，所述量子比特具有与器件物理结构相关的噪声特征，包括但不限于 $T_1$ 弛豫时间、$T_2/T_2^*$ 退相干时间、单比特门误差、双比特门误差、读出误差和串扰系数。参数化电路执行模块用于执行变分量子电路、校准电路、随机基准测试电路及其组合，使量子计算结果和设备噪声特征能够同时反映于测量输出中。

机器学习表征模块用于基于测量统计、校准遥测和运行元数据提取设备级噪声参数，区分相干信号、自然漂移与异常噪声。误差映射与指纹生成模块将高维连续噪声向量归一化、量化和哈希映射，生成硬件指纹和设备绑定签名。密钥生成模块基于指纹结果经哈希函数和密钥派生函数得到后量子兼容的密钥种子或会话密钥。篡改检测模块持续比较当前噪声画像与注册基线的偏离程度，在超阈值时触发告警、密钥失效、重新校准和/或重新注册流程。部署接口模块将上述能力以云 API、验证服务、中间件适配器、控制平面固件或 qHSM 方式对外提供。

\subsection{核心方法步骤}

\subsubsection{1. 噪声特征采集与数学表征}
本发明优选采集以下设备噪声参数：纵向弛豫时间 $T_1$、横向退相干时间 $T_2$ 或 $T_2^*$、单比特门平均失真率 $\epsilon_1$、双比特门平均失真率 $\epsilon_2$、读出误差 $\epsilon_r$ 以及串扰系数 $c_{ij}$。

\subsubsection{2. 硬件指纹构造}
本发明将量子设备噪声参数组织为硬件指纹向量：
\[
\mathbf{f} = [T_1^{(1)}, T_2^{(1)}, \ldots, T_1^{(n)}, T_2^{(n)}, \epsilon_1^{(1)}, \ldots, \epsilon_2^{(i,j)}, \ldots, c_{12}, c_{13}, \ldots, \epsilon_r^{(1)}, \ldots ]
\]

对于具有 $n$ 个量子比特、$m$ 个双比特门配置的器件，其指纹维度可包含 $2n$ 个弛豫/退相干参数、$n$ 个单比特误差参数、$n$ 个读出误差参数、$m$ 个双比特门误差参数以及 $n(n-1)/2$ 个串扰参数。

\subsubsection{3. 指纹到密钥的映射}
由于量子噪声值为连续实数，本发明优选采用“归一化 + 自适应量化 + 模糊提取 + 哈希/KDF”流水线实现稳定密钥导出。可采用下式表示密钥导出形式：
\[
K = \mathrm{KDF}(\mathrm{Hash}(\mathbf{f}) \| \mathrm{salt} \| \mathrm{context})
\]
其中，$\mathbf{f}$ 为指纹向量，$\mathrm{salt}$ 为随机或策略上下文，$\mathrm{context}$ 为设备身份、租户策略、用途标识等。

\subsubsection{4. 篡改检测机制}
设注册阶段得到的基线指纹为 $\mathbf{f}_{baseline}$，当前采样得到的指纹为 $\mathbf{f}_{current}$，则可以 Mahalanobis 距离为例：
\[
d = \sqrt{(\mathbf{f}_{current}-\mathbf{f}_{baseline})^T \Sigma^{-1}(\mathbf{f}_{current}-\mathbf{f}_{baseline})}
\]
其中，$\Sigma$ 表示自然漂移协方差矩阵。当 $d$ 超过阈值 $\tau$ 时，系统触发篡改告警、使当前派生密钥失效、阻断后续证明或作业调度、记录审计日志并通知上层系统，以及要求重新注册或重新校准中的一项或多项动作。

\subsubsection{5. 工程化部署结构}
本发明优选划分如下可部署子系统：控制面/设备侧采集器、特征提取服务、指纹与密钥服务、验证与策略引擎以及企业集成接口。在优选实施方式中，还可设置量子中间件抽象层，使不同厂商量子云、控制器和校准总线的差异被适配器封装，上层证明与指纹逻辑无需随平台变化而整体重写。

\subsection{实施方式与实验示例}

采用 27 比特超导量子处理器作为测试对象，对全部量子比特按小时执行表征任务，在 30 天内完成 720 个表征周期，总电路数超过 500 万，单电路 shots 为 4096。基线测得参数范围如下：$T_1$ 为 72--198~$\mu$s，均值 142~$\mu$s；$T_2$ 为 41--167~$\mu$s，均值 98~$\mu$s；单比特误差为 0.02\%--0.15\%，均值 0.06\%；双比特误差为 0.4\%--1.8\%，均值 0.9\%；读出误差为 0.5\%--2.1\%，均值 1.2\%。

选取 5 台不同量子处理器进行比较，不同设备之间的欧氏距离均值约为 0.44，Mahalanobis 距离均值约为 9.2；同一设备在 30 天内不同时刻的 Mahalanobis 距离约为 0.8--2.8。根据实验结果，可设置“同设备阈值”为 $d < 4.0$。实验统计显示：假接受率小于 0.001\%，假拒绝率小于 0.1\%，等错误率约为 0.02\%，原始指纹熵约为 187 bits。

将噪声特征按参数类别映射为不同熵贡献，经条件化后，可生成约 128 bits 安全强度的密钥。典型导出时间包括：指纹提取 8~ms、量化 2~ms、SHA3-256 哈希 1~ms、HKDF 扩展 3~ms，总计约 14~ms。对派生结果执行 NIST SP 800-22 测试，频数、块频数、游程、傅里叶变换、近似熵等项目均通过。

\section{本申请提案的关键点和欲保护点}

本申请提案主要包含以下关键创新点和保护点：

(1) \textbf{保护一种基于量子处理器固有噪声参数构建设备身份的方法和系统}：以 $T_1$、$T_2$、门误差、串扰和读出误差等物理参数作为设备原生身份特征，区别于基于经典制造偏差或单纯随机数输出的现有方案。

(2) \textbf{保护一种从连续量子噪声特征生成可重复硬件指纹的映射机制}：包括对噪声参数进行采集、归一化、量化、压缩、哈希或编码处理，以得到可重复、可区分且难以伪造的设备指纹。

(3) \textbf{保护一种利用量子噪声指纹导出设备绑定密钥或证明材料的机制}：通过模糊提取、哈希函数和密钥派生函数，将噪声指纹进一步转化为设备绑定密钥、证明令牌或认证结果。

(4) \textbf{保护一种区分自然漂移与恶意篡改的运行期检测机制}：利用 Mahalanobis 距离或其他多维统计距离衡量当前噪声画像与注册基线的偏离程度，并在超阈值时触发告警、密钥失效、证明拒绝、重新校准和重新注册中的一项或多项处理。

(5) \textbf{保护一种面向量子云、控制平面、中间件和 qHSM 的工程化部署架构}：包括设备侧采集器、特征提取服务、指纹/密钥服务、验证服务、审计与策略引擎以及适配不同量子平台的抽象接口。

(6) \textbf{保护上述技术方案在系统实现中的整体协同关系}：使设备身份构建、指纹生成、密钥导出、篡改检测和部署接口构成完整可实施的技术方案。

\section{与第三条中最接近的现有技术相比，本申请提案有何技术优点}

相比于现有技术（如经典 PUF、量子随机数生成方案、一般量子设备证明方案或单纯量子噪声表征方案），本方案具有显著优势：

(1) \textbf{有助于形成可验证的设备绑定能力}：本发明以量子处理器物理层噪声参数作为身份根，不依赖可拷贝的软件证书、通用密钥文件或单纯逻辑标识。

(2) \textbf{兼顾区分度与可重构性}：通过对高维噪声向量进行归一化、量化和模糊提取，可以在同一设备自然漂移范围内提高密钥重构一致性，同时使不同设备间保持显著区分度。

(3) \textbf{支持运行期持续证明和篡改告警}：不仅完成一次性识别，还可在运行期间持续比较当前噪声画像与历史基线之间的偏离，实现对探针插入、电磁扰动、温漂异常等事件的检测与响应。

(4) \textbf{具备工程可部署性}：本发明不仅给出理论模型，而且提供了面向量子云、控制平面、中间件、验证服务和 qHSM 的产品化落地路径。

\section{其他有助于理解本申请提案的技术资料}

(1) Magesan, Gambetta, Emerson, ``Scalable and Robust Randomized Benchmarking of Quantum Processes'', PRL, 2011.

(2) Nielsen \& Chuang, ``Prescription for Experimental Determination of the Dynamics of a Quantum Black Box'', J. Mod. Opt., 1997.

(3) Arapinis et al., ``Quantum Physical Unclonable Functions'', arXiv:1905.02550, 2019.

(4) Herrero-Collantes \& Garcia-Escartin, ``Quantum Random Number Generators'', Rev. Mod. Phys., 2017.

\section{本申请提案的侵权证据可获得性}

本提案的侵权行为具有较高的可检测性：

(1) \textbf{接口与服务证据}：若第三方对外提供量子硬件证明、设备绑定密钥或基于量子噪声的认证接口，可通过 API 文档、产品白皮书、采购技术规范和测试环境输出识别其是否落入本方案技术特征。

(2) \textbf{系统部署证据}：若公开产品支持量子云、中间件或 qHSM 形态下的噪声画像采集、指纹映射和篡改联动机制，则可从产品架构说明、审计日志、配置项、运行策略和合规材料中获得证据。

(3) \textbf{运行行为证据}：在具备测试权限的情况下，可通过多轮注册、证明、异常注入和密钥更新流程观察其是否存在本发明特有的“噪声采集—指纹生成—密钥导出—偏离检测—自动响应”闭环。

\end{document}
